\chapter{Estado común, árbol de prolongaciones y medida genealógica}
\label{chap:continuo-estado-arbol-medida-comun}

\begin{thesisbox}[title={Un solo estado enriquecido antes de toda discriminación interna}]
El objeto de este capítulo es el estado finito completo producido por
\(\APP\), orientado por TRIT y transportado por \(\TPK\).  A partir de él se
construyen el grupoide de simetrías reversibles, el árbol de todas las
prolongaciones admisibles, la multisección de historias completas y su medida
canónica.  No se introduce ningún dominio especializado: el dominio es el
conjunto total de estados que efectivamente genera el mismo mecanismo
\(\APP\)--TRIT--\(\TPK\).  Tampoco se selecciona una hoja, una continuación o
una lectura posterior.
\end{thesisbox}

\section{Datos finitos producidos por APP}
\label{sec:continuo-comun-app-finito}

Para cada profundidad \(k\geq0\), denotamos por \(\mathsf A_k\) el conjunto
finito de prefijos celulares producidos por APP y por
\begin{equation}
\label{eq:continuo-comun-app-truncamiento}
 \rho^{\APP}_{k+1,k}:\mathsf A_{k+1}\longrightarrow\mathsf A_k
\end{equation}
el borrado del último refinamiento.  Este borrado conserva el prefijo ya
construido.  La célula neutra de APP da una prolongación
\begin{equation}
\label{eq:continuo-comun-app-degeneracion}
 \eta^{\APP}_k:\mathsf A_k\longrightarrow\mathsf A_{k+1},
 \qquad
 \rho^{\APP}_{k+1,k}\eta^{\APP}_k=I_{\mathsf A_k}.
\end{equation}
Escribimos
\begin{equation}
\label{eq:continuo-comun-emision-neutra-app}
 \varepsilon^0_k(a):a\longrightarrow\eta^{\APP}_k(a)
\end{equation}
para la emisión que realiza esa prolongación.  Así se mantienen separados su
destino, \(\eta^{\APP}_k(a)\), y la operación que lo produce,
\(\varepsilon^0_k(a)\).
No se interpreta \(\eta^{\APP}_k\) como ausencia de historia: registra un
refinamiento de longitud uno con emisión neutra y conserva materialmente el
estado anterior.

Cada \(a\in\mathsf A_k\) porta simultáneamente las dos hojas de APP,
\begin{equation}
\label{eq:continuo-comun-dos-hojas}
 \operatorname{Sh}_k(a)
 =\bigl(\Sigma_k(a),\Pi_k(a)\bigr)
 \in\mathsf H_k^{\Sigma}\times\mathsf H_k^{\Pi},
\end{equation}
y no un elemento elegido de la pareja.  Para una emisión elemental
\(\epsilon:a\to a'\), APP produce además dos divisiones euclídeas tipadas,
\begin{equation}
\label{eq:continuo-comun-residuo-cociente-app}
 \delta^{\diamond}(\epsilon)
 =\operatorname{res}^{\diamond}(\epsilon)
  +9\operatorname{quo}^{\diamond}(\epsilon),
 \qquad
 \operatorname{res}^{\diamond}(\epsilon)\in\{0,\ldots,8\},
 \quad
 \operatorname{quo}^{\diamond}(\epsilon)\in\mathbb Z,
\end{equation}
donde \(\diamond\in\{\Sigma,\Pi\}\).  La etiqueta \(\diamond\) sólo señala
qué hoja se está leyendo; ambas divisiones permanecen en el mismo registro.

La fuente y el destino celulares de \(\epsilon\) determinan su frontera
orientada
\begin{equation}
\label{eq:continuo-comun-frontera-elemental}
 \partial\epsilon:=t(\epsilon)-s(\epsilon).
\end{equation}
La emisión neutra satisface
\begin{equation}
\label{eq:continuo-comun-neutral-app-datos}
 \operatorname{res}^{\diamond}(\varepsilon^0_k(a))=0,
 \qquad
 \operatorname{quo}^{\diamond}(\varepsilon^0_k(a))=0,
 \qquad
 \partial\varepsilon^0_k(a)=0.
\end{equation}

\begin{definition}[Grafo dirigido de emisiones APP]
\label{def:continuo-comun-grafo-app}
Sea \(\mathsf E_k(a)\) el conjunto finito de emisiones APP de longitud uno
que parten de \(a\in\mathsf A_k\).  Conservamos las emisiones con
multiplicidad: dos operaciones con el mismo destino pero con distinto
residuo, cociente o procedencia son elementos diferentes de
\(\mathsf E_k(a)\).  Por \eqref{eq:continuo-comun-app-degeneracion},
\begin{equation}
\label{eq:continuo-comun-out-no-vacio-app}
 \varepsilon^0_k(a)\in\mathsf E_k(a),
 \qquad
 1\leq \#\mathsf E_k(a)<\infty.
\end{equation}
\end{definition}

\section{Levantamiento TRIT: orientación y acarreo genealógico producidos}
\label{sec:continuo-comun-trit-lift}

La incidencia firmada de una emisión APP es un entero
\(\Delta(\epsilon)\).  El representante equilibrado módulo tres da de forma
única
\begin{equation}
\label{eq:continuo-comun-trit-euclid}
 \Delta(\epsilon)
 =\operatorname{or}(\epsilon)+3\kappa(\epsilon),
 \qquad
 \operatorname{or}(\epsilon)\in\{-1,0,+1\},
 \quad
 \kappa(\epsilon)\in\mathbb Z.
\end{equation}
Así, la orientación y el carry no son condiciones impuestas a una historia:
son funciones totales del dato APP.

Para hacer visible la composición, sea
\(\mathsf C_{\epsilon}:\mathsf M_{s(\epsilon)}\to
\mathsf M_{t(\epsilon)}\) el transporte de memoria inducido por la emisión.
El carry nativo de una composición cumple
\begin{equation}
\label{eq:continuo-comun-cociclo-carry}
 \kappa(\epsilon_2\epsilon_1)
 =\kappa(\epsilon_2)
  +\mathsf C_{\epsilon_2}\kappa(\epsilon_1),
\end{equation}
siempre que \(t(\epsilon_1)=s(\epsilon_2)\).  Para la identidad y la emisión
neutra,
\begin{equation}
\label{eq:continuo-comun-carry-unidad}
 \kappa(I_a)=0,
 \qquad
 \kappa(\varepsilon^0_k(a))=0,
 \qquad
 \mathsf C_{I_a}=I_{\mathsf M_a}.
\end{equation}

\begin{lemma}[Asociatividad del ledger de carry]
\label{lem:continuo-comun-asociatividad-carry}
Para tres emisiones componibles, la evaluación de
\(\kappa(\epsilon_3\epsilon_2\epsilon_1)\) no depende del parentizado.
\end{lemma}

\begin{proof}
Aplicando dos veces \eqref{eq:continuo-comun-cociclo-carry} al parentizado
izquierdo se obtiene
\[
 \kappa(\epsilon_3)
 +\mathsf C_{\epsilon_3}\kappa(\epsilon_2)
 +\mathsf C_{\epsilon_3}\mathsf C_{\epsilon_2}
  \kappa(\epsilon_1).
\]
El parentizado derecho produce la misma expresión porque el transporte de
memoria es functorial:
\(\mathsf C_{\epsilon_3\epsilon_2}
=\mathsf C_{\epsilon_3}\mathsf C_{\epsilon_2}\).
\end{proof}

La reversión orientada de una emisión reversible se denota
\(\overline\epsilon\).  TRIT actúa por
\begin{equation}
\label{eq:continuo-comun-reversion-trit}
 \operatorname{or}(\overline\epsilon)
 =-\operatorname{or}(\epsilon),
 \qquad
 \partial\overline\epsilon=-\partial\epsilon,
 \qquad
 \mathsf C_{\overline\epsilon}=\mathsf C_\epsilon^{-1}.
\end{equation}
El carry inverso queda determinado, no elegido, por
\begin{equation}
\label{eq:continuo-comun-carry-inverso}
 \kappa(\overline\epsilon)
 =-\mathsf C_\epsilon^{-1}\kappa(\epsilon),
\end{equation}
pues \(\kappa(\overline\epsilon\epsilon)=0\).

\section{Transporte TPK de las \(2\) hojas y de la memoria}
\label{sec:continuo-comun-tpk-transporte}

La cronolog\'ia de TPK forma parte del tipo de cada emisi\'on y no se
reconstruye a posteriori contando letras. Recuperamos aqu\'i la acci\'on del
od\'ometro de \cref{def:odometro-9-12,prop:proyeccion-fase-longitud}. Sea
\begin{equation}
\label{eq:continuo-comun-ciclo-fase-tpk}
 C_9:=\mathbb Z/9\mathbb Z,\qquad
 \widetilde{\mathcal O}_{9}:=\mathbb Z_{\geq0}\times C_9,\qquad
 \mathcal O_{9,12}:=\mathbb Z/12\mathbb Z\times C_9.
\end{equation}
El estado desenrollado de fase y memoria es
\(\mathsf Q_k=(w_k,r_k)\in\widetilde{\mathcal O}_{9}\); su reducción
\((w,r)\mapsto([w]_{12},r)\) recupera
\(\mathcal O_{9,12}\). La semilla canónica del reloj es
\begin{equation}
\label{eq:continuo-comun-semilla-fase-tpk}
 \mathsf Q_0(a_0):=(0,[0]_9)\qquad(a_0\in\mathsf A_0);
\end{equation}
las demás elecciones de origen son sus traslaciones coherentes por \(C_9\)
y no cambian la ley de actualización. Para
\(\bar r\in\{0,\ldots,8\}\), ponemos
\begin{equation}
\label{eq:continuo-comun-carry-fase-tpk}
 \chi_9(r):=\left\lfloor\frac{\bar r+1}{9}\right\rfloor.
\end{equation}
Toda emisión \(\epsilon\) que lleva profundidad \(k\) a profundidad \(k+1\)
actúa en el odómetro por
\begin{equation}
\label{eq:continuo-comun-accion-fase-tpk}
 \sigma_9(w,r):=(w+\chi_9(r),r+[1]_9),\qquad
 \mathsf Q_{k+1}(\epsilon_*\widehat x_k)
 =\sigma_9\mathsf Q_k(\widehat x_k).
\end{equation}
Su componente de fase es, por definición tipada,
\begin{equation}
\label{eq:continuo-comun-fase-emision-tpk}
 \operatorname{ph}(\epsilon):r_k\longmapsto r_k+[1]_9,
 \qquad
 \operatorname{ph}(\epsilon_n\cdots\epsilon_1):
 r_k\longmapsto r_k+[n]_9.
\end{equation}
Escribimos \(g_k:=r_k\in C_9\) para la fase publicada y
\(q^{(9)}_k:=w_k\) para la memoria vertical, y
\(\beta_k:=[w_k]_{12}\) para su reducción dodecafásica. La iteración de
\eqref{eq:continuo-comun-accion-fase-tpk} es
\begin{equation}
\label{eq:continuo-comun-iteracion-fase-tpk}
 \sigma_9^n(w,r)=
 \left(w+\left\lfloor\frac{\bar r+n}{9}\right\rfloor,
       r+[n]_9\right).
\end{equation}
En efecto, los sumandos
\(\chi_9(r),\chi_9(r+[1]_9),\ldots,\chi_9(r+[n-1]_9)\)
cuentan exactamente los cruces del representante \(8\) al representante
\(0\), que son
\(\lfloor(\bar r+n)/9\rfloor\). En particular,
\begin{equation}
\label{eq:continuo-comun-retorno-fase-avance-memoria}
 \sigma_9^9(w,r)=(w+1,r),\qquad
 g_{k+9}=g_k,\qquad q^{(9)}_{k+9}=q^{(9)}_k+1.
\end{equation}
Incluso la emisión neutra de APP realiza
\(g_k\mapsto g_{k+1}\): es neutra en residuo, cociente, carry y frontera,
pero registra que ha ocurrido una extensión. La palabra de fase
\begin{equation}
\label{eq:continuo-comun-palabra-fase-tpk}
 \boldsymbol\tau_k^{(N)}:=(g_k,g_{k+1},\ldots,g_N),
 \qquad g_{j+1}=g_j+[1]_9,
\end{equation}
y la memoria \(q^{(9)}\) tipan por objetos distintos el retorno de fase y
el cambio de estado.

\subsection{Renovación simultánea de la fibra completa}
\label{sec:continuo-comun-kph}

El odómetro anterior conserva la cronología de cada extensión. La misma
vuelta nonádica posee además la realización Markov interna, ya construida
sobre el catálogo APP--TPK, que renueva simultáneamente las dos coordenadas
de la fibra completa. Es esencial conservar ambas realizaciones y no
identificar una arista genérica del árbol con una emisión del catálogo sin
un mapa de tipado.

\begin{proposition}[Realización Markov interna de la conexión nonádica]
\label{prop:continuo-comun-kph-heredada}
\label{pII-full:thm:seleccion-novena-fase}
Sean el dominio de semillas APP, su emisión observable y sus multiplicidades
\[
 T_{\min}:\Omega_{\rm seed}\twoheadrightarrow\mathcal U_6,
 \qquad
 m(u):=\#T_{\min}^{-1}(u),
\]
y la factorización producida por el TPK
\[
 \mathcal U_6\simeq\mathcal S_+\times\mathcal S_\times,
 \qquad |\mathcal S_+|=18,\qquad |\mathcal S_\times|=26.
\]
Sobre funciones en \(\mathcal U_6\), las dos esperanzas condicionales actúan
por
\[
 (E_+f)(s,e)=\frac1{18}\sum_{s'\in\mathcal S_+}f(s',e),
 \qquad
 (E_\times f)(s,e)=\frac1{26}
 \sum_{e'\in\mathcal S_\times}f(s,e').
\]
Para
\[
 \tau=(+1,-1,0,+1,-1,0,+1,-1,0),
 \qquad
 P_{+1}=E_+,\quad P_{-1}=E_\times,\quad P_0=I,
\]
la transferencia de catálogo y su elevación a todas las semillas son
\begin{align}
 \mathcal K_{\rm ph}\bigl((u,r),(v,r+1)\bigr)
 &=P_{\tau(r)}(u,v),
 \label{eq:continuo-comun-kph-catalogo}\\
 L_0(\omega,\omega')&=\mathbf1_{\{\omega=\omega'\}},\nonumber\\
 L_+(\omega,\omega')
 &=\frac{E_+(T_{\min}\omega,T_{\min}\omega')}
         {m(T_{\min}\omega')},\nonumber\\
 L_\times(\omega,\omega')
 &=\frac{E_\times(T_{\min}\omega,T_{\min}\omega')}
         {m(T_{\min}\omega')},\nonumber\\
 \widehat{\mathcal K}_{\rm ph}
 \bigl((\omega,r),(\omega',r+1)\bigr)
 &=L_{\tau(r)}(\omega,\omega').
 \label{eq:continuo-comun-kph-elevada}
\end{align}
Para \(u=T_{\min}\omega\), se cumple fase a fase
\begin{equation}
\label{eq:continuo-comun-kph-lumpability}
 \sum_{\omega':\,T_{\min}\omega'=v}
 \widehat{\mathcal K}_{\rm ph}
 \bigl((\omega,r),(\omega',r+1)\bigr)
 =
 \mathcal K_{\rm ph}\bigl((u,r),(v,r+1)\bigr).
\end{equation}
En particular,
\begin{equation}
\label{eq:continuo-comun-kph-novena-potencia}
 \boxed{\;
 \mathcal K_{\rm ph}^{\,9}\big|_r=E_+E_\times,
 \qquad
 (E_+E_\times)(u,v)=\frac1{468}.
 \;}
\end{equation}
\end{proposition}

\begin{proof}
En una fase \(+\) o \(\times\), sumar
\eqref{eq:continuo-comun-kph-elevada} sobre
\(T_{\min}^{-1}(v)\) cancela exactamente el denominador \(m(v)\); en una
fase cero, \(L_0\) proyecta a la identidad. Esto demuestra
\eqref{eq:continuo-comun-kph-lumpability}. Los operadores \(E_+\),
\(E_\times\) e \(I\) son idempotentes y conmutan, y cada uno aparece tres
veces en la palabra nonádica. Por tanto
\[
 \prod_{r\in C_9}P_{\tau(r)}
 =E_+^3E_\times^3I^3=E_+E_\times.
\]
La composición renueva primero una coordenada de cardinal \(18\) y después
la otra de cardinal \(26\), de modo que cada destino recibe
\(1/(18\cdot26)=1/468\).
\end{proof}

\begin{remark}[Dos realizaciones internas, sin identificación ficticia]
\label{rem:continuo-comun-kph-no-identificacion-aristas}
La transferencia \(\mathcal K_{\rm ph}\) es la realización global del
catálogo completo y \(\sigma_9\) es la acción cronológica sobre cada estado
del árbol. No se sustituye una por otra. Una igualdad arista a arista
requeriría un mapa adicional
\[
 \operatorname{Cat}_{k,x}:\operatorname{Out}(x)\longrightarrow\mathcal U_6
\]
que verificase la correspondiente identidad de suma sobre fibras. Ese mapa
no se usa en este volumen. La monodromía completa
\eqref{eq:continuo-comun-kph-novena-potencia} procede del catálogo
APP--TPK, mientras la genealogía de cada historia permanece en el árbol y
su ledger.
\end{remark}

TPK asigna a cada emisión admisible un transporte tipado de la pareja de
hojas,
\begin{equation}
\label{eq:continuo-comun-transporte-hojas}
 U_\epsilon:
 \mathsf H_{s(\epsilon)}^{\Sigma}\times
 \mathsf H_{s(\epsilon)}^{\Pi}
 \longrightarrow
 \mathsf H_{t(\epsilon)}^{\Sigma}\times
 \mathsf H_{t(\epsilon)}^{\Pi},
\end{equation}
con
\begin{equation}
\label{eq:continuo-comun-transporte-hojas-funtorial}
 U_{I_a}=I,
 \qquad
 U_{\epsilon_2\epsilon_1}=U_{\epsilon_2}U_{\epsilon_1}.
\end{equation}
El codominio de \(U_\epsilon\) sigue siendo la pareja completa.  Ninguna
operación de este capítulo sustituye
\((\Sigma,\Pi)\) por una de sus coordenadas.

Sea \(m\in\mathsf M_{s(\epsilon)}\) una memoria acumulada.  La acción TPK
sobre ella es la transformación afín
\begin{equation}
\label{eq:continuo-comun-accion-memoria}
 \mathsf U_\epsilon(m)
 :=\mathsf C_\epsilon m+\kappa(\epsilon).
\end{equation}

\begin{lemma}[Acción TPK sobre memoria]
\label{lem:continuo-comun-accion-memoria}
Las transformaciones \(\mathsf U_\epsilon\) satisfacen
\begin{equation}
\label{eq:continuo-comun-accion-memoria-composicion}
 \mathsf U_{I_a}=I,
 \qquad
 \mathsf U_{\epsilon_2\epsilon_1}
 =\mathsf U_{\epsilon_2}\mathsf U_{\epsilon_1}.
\end{equation}
\end{lemma}

\begin{proof}
La identidad se sigue de \eqref{eq:continuo-comun-carry-unidad}.  Para la
composición,
\begin{align*}
 \mathsf U_{\epsilon_2}\mathsf U_{\epsilon_1}(m)
 &=\mathsf C_{\epsilon_2}
   \bigl(\mathsf C_{\epsilon_1}m+\kappa(\epsilon_1)\bigr)
   +\kappa(\epsilon_2)\\
 &=\mathsf C_{\epsilon_2\epsilon_1}m
   +\kappa(\epsilon_2)
   +\mathsf C_{\epsilon_2}\kappa(\epsilon_1)\\
 &=\mathsf C_{\epsilon_2\epsilon_1}m
   +\kappa(\epsilon_2\epsilon_1),
\end{align*}
donde la última igualdad es
\eqref{eq:continuo-comun-cociclo-carry}.
\end{proof}

Una ruta finita es una palabra componible
\begin{equation}
\label{eq:continuo-comun-ruta}
 \gamma=\epsilon_1\epsilon_2\cdots\epsilon_k.
\end{equation}
Su frontera, transporte, carry y memoria son
\begin{align}
 \partial\gamma
 &:=\sum_{j=1}^{k}\partial\epsilon_j,
 \label{eq:continuo-comun-frontera-ruta}\\
 U_\gamma
 &:=U_{\epsilon_k}\cdots U_{\epsilon_1},
 \label{eq:continuo-comun-transporte-ruta}\\
 \kappa(\gamma)
 &:=\kappa(\epsilon_k)
   +\sum_{j=1}^{k-1}
    \mathsf C_{\epsilon_k}\cdots\mathsf C_{\epsilon_{j+1}}
    \kappa(\epsilon_j),
 \label{eq:continuo-comun-carry-ruta}\\
 m(\gamma;m_0)
 &:=\mathsf U_{\epsilon_k}\cdots
       \mathsf U_{\epsilon_1}(m_0).
 \label{eq:continuo-comun-memoria-ruta}
\end{align}
En \eqref{eq:continuo-comun-frontera-ruta}, los extremos interiores se
cancelan y queda
\begin{equation}
\label{eq:continuo-comun-frontera-telescopica}
 \partial\gamma=t(\epsilon_k)-s(\epsilon_1).
\end{equation}

\section{El estado común y su dominio total}
\label{sec:continuo-comun-estado-total}

\begin{definition}[Ledger elemental]
\label{def:continuo-comun-ledger-elemental}
Para una emisión que aparece entre dos prefijos enriquecidos consecutivos,
denotamos por
\(\mathsf Q^-(\epsilon)=(w^-(\epsilon),r^-(\epsilon))\) el odómetro del
prefijo fuente y por
\begin{equation}
\label{eq:continuo-comun-ledger-fase-extremos}
 \mathsf Q^+(\epsilon):=\sigma_9\mathsf Q^-(\epsilon)
 =(w^-(\epsilon)+\chi_9(r^-(\epsilon)),r^-(\epsilon)+[1]_9)
\end{equation}
el del prefijo destino. Así no se aplica el odómetro a una célula APP
desnuda. El registro de la emisión es
\begin{equation}
\label{eq:continuo-comun-ledger-elemental}
 \lambda(\epsilon):=\bigl(
 s(\epsilon),t(\epsilon),
 \operatorname{res}^{\Sigma}(\epsilon),
 \operatorname{quo}^{\Sigma}(\epsilon),
 \operatorname{res}^{\Pi}(\epsilon),
 \operatorname{quo}^{\Pi}(\epsilon),
 \operatorname{or}(\epsilon),
 \kappa(\epsilon),
 \partial\epsilon,
 \mathsf Q^-(\epsilon),\mathsf Q^+(\epsilon)\bigr).
\end{equation}
El ledger de \(\gamma\) es la palabra ordenada
\begin{equation}
\label{eq:continuo-comun-ledger-ruta}
 \operatorname{Led}(\gamma)
 :=\bigl(\lambda(\epsilon_1),\ldots,
          \lambda(\epsilon_k)\bigr).
\end{equation}
Conservar esta palabra, y no sólo su suma, distingue dos historias con la
misma salida terminal.
\end{definition}

\subsection{Lector completo del estado TPK}
\label{sec:continuo-comun-lector-estado-tpk}

La profundidad \(K\) de bloques TPK y la profundidad \(k\) del árbol de
extensiones son tipos distintos. En esta subsección \(K\) denota
exclusivamente la profundidad del estado TPK ya generado por el mecanismo
operatorio; no se identifica con \(k\) ni con una arista genérica.

\begin{definition}[Lector compacto de memoria nonádica]
\label{def:continuo-comun-lector-compacto-nonadico}
El estado enriquecido TPK de profundidad \(K\) es
\begin{equation}
\label{eq:continuo-comun-estado-tpk-completo}
 v_K=
 \bigl(
 B_K,\boldsymbol\delta_K,\boldsymbol x_K,
 \boldsymbol{\mathcal C}_K,\Sigma_K,
 \omega_K,\ell_K,g(K),\Lambda_K
 \bigr),
 \qquad
 B_K=(u_K^\pi,u_K^e,u_K^\varphi).
\end{equation}
Su lector compacto no añade coordenadas: las publica desde ese mismo estado,
\begin{equation}
\label{eq:continuo-comun-lector-compacto-nonadico}
 \mathfrak r_9(v_K)
 :=\bigl(B_K,C_K,p_K,c_K,\Sigma_K,W_K,g(K)\bigr),
\end{equation}
donde
\begin{equation}
\label{eq:continuo-comun-lector-compacto-componentes}
\begin{aligned}
 C_K&:=\boldsymbol{\mathcal C}_K
      =\bigl(C_K^\chi\bigr)_{\chi\in\{\pi,e,\varphi\}},\\
 C_K^\chi&:=(N_K^\chi,L_K^\chi,r_K^\chi,
             D_K^\chi,A_K^\chi,B_K^\chi),\\
 p_K&:=(B_0,\ldots,B_K),\\
 c_K&:=(\boldsymbol\delta_K,\boldsymbol x_K,
       \kappa(\gamma^{\rm TPK}_{0:K})),
\end{aligned}
\end{equation}
y la sombra orientada de Witt se conserva como
\begin{equation}
\label{eq:continuo-comun-sombra-witt}
 W_K:=
 \bigl((u_K^\chi,u_K^\chi A_W)\bigr)_
       {\chi\in\{\pi,e,\varphi\}},
 \qquad
 A_W=
 \begin{pmatrix}
 0&1&1&1&1&1\\
 1&0&1&2&2&1\\
 1&1&0&1&2&2\\
 1&2&1&0&1&2\\
 1&2&2&1&0&1\\
 1&1&2&2&1&0
 \end{pmatrix},
 \qquad A_W^2=-I_6.
\end{equation}
La ruta \(\gamma^{\rm TPK}_{0:K}\) es la palabra de extensiones codificada
por \(\Lambda_K\); su prefijo de longitud \(j\) es el que produce
\(v_j\). La firma es la salida de la primitiva interna explícita
\begin{equation}
\label{eq:continuo-comun-primitiva-firma-tpk}
 \Sigma_B:
 \mathsf B_{\rm TPK}\times\mathsf C_{\rm TPK}
 \times\mathbb Z[\mathsf A]\times\Omega_{12}\times\mathsf L_{\rm TPK}
 \longrightarrow\mathsf S_{\rm TPK},
 \quad
 (B,c,\partial\gamma,\omega,\ell)\longmapsto
 \Sigma_B(B,c,\partial\gamma,\omega,\ell),
\end{equation}
con \(\Sigma_B(g\cdot-)=g\cdot\Sigma_B(-)\) para cada simetría TPK. Su
estructura de partida explícita es
\((B_K,c_K,\partial\gamma^{\rm TPK}_{0:K},\omega_K,\ell_K)\). No se la reemplaza por una
firma exterior ni se infiere retrospectivamente desde un lector numérico.
\end{definition}

La acción del paso \(K\mapsto K+1\) se escribe sin omitir el estado profundo.
Para cada \(\chi\in\{\pi,e,\varphi\}\),
\begin{equation}
\label{eq:continuo-comun-update-hensel-tpk}
 y_K^\chi=x_K^\chi A_H,\qquad
 u_{K+1}^\chi=y_K^\chi\bmod3,\qquad
 x_{K+1}^\chi=\frac{y_K^\chi-u_{K+1}^\chi}{3},
\end{equation}
con
\[
 A_H=
 \begin{pmatrix}
 2&2&2&1&2&1\\
 2&1&2&2&1&1\\
 1&0&0&1&0&2\\
 0&0&1&1&1&0\\
 2&2&2&0&2&1\\
 2&1&2&1&0&0
 \end{pmatrix},
 \qquad \det A_H=1.
\]
Si
\[
 \operatorname{ind}_3(u_1,\ldots,u_6)
 :=\sum_{j=1}^{6}\overline u_j\,3^{6-j}
 \in\{0,\ldots,728\},
\]
la actualización del prefijo y de sus \(729\) subcilindros es
\begin{align}
 N_{K+1}^\chi
 &=729N_K^\chi+\operatorname{ind}_3(u_{K+1}^\chi),
 &L_{K+1}^\chi&=L_K^\chi+6,
 \label{eq:continuo-comun-update-cilindro-indice}\\
 I_{K,u}^\chi
 &:=
 \left[
  \frac{729N_K^\chi+u}{3^{L_K^\chi+6}},
  \frac{729N_K^\chi+u+1}{3^{L_K^\chi+6}}
 \right),
 &0&\leq u<729.
 \label{eq:continuo-comun-update-cilindro-intervalo}
\end{align}
El cilindro decimal y el acarreo asociado se actualizan por los dos casos
enteros del mismo operador \(\mathcal C_{3,1000}\). Si no se publica una
nueva tríada decimal,
\begin{equation}
\label{eq:continuo-comun-update-cilindro-sin-triada}
 r_{K+1}^\chi=r_K^\chi,\qquad D_{K+1}^\chi=D_K^\chi,\qquad
 A_{K+1}^\chi=729A_K^\chi+
 \operatorname{ind}_3(u_{K+1}^\chi)1000^{r_K^\chi},\qquad
 B_{K+1}^\chi=A_{K+1}^\chi+1000^{r_K^\chi}.
\end{equation}
La primitiva \(\mathcal C_{3,1000}\) produce el dato opción
\[
 d_{K+1,\chi}^{\rm dec}\in
 \{\bot\}\sqcup\{0,\ldots,999\};
\]
en el primer caso vale \(d_{K+1,\chi}^{\rm dec}=\bot\). Si aparece una
tríada, escribimos
\(d_{K+1,\chi}^{\rm dec}=\delta_{K+1}^\chi\in\{0,\ldots,999\}\) y
\begin{equation}
\label{eq:continuo-comun-update-cilindro-con-triada}
 \begin{aligned}
 r_{K+1}^\chi&=r_K^\chi+1,\qquad
 D_{K+1}^\chi=1000D_K^\chi+\delta_{K+1}^\chi,\\
 A_{K+1}^\chi&=729000A_K^\chi
 +\operatorname{ind}_3(u_{K+1}^\chi)1000^{r_K^\chi+1}
 -729\delta_{K+1}^\chi3^{L_K^\chi},\\
 B_{K+1}^\chi&=A_{K+1}^\chi+1000^{r_K^\chi+1}.
 \end{aligned}
\end{equation}
El registro \(\boldsymbol\delta_{K+1}\) añade esas tríadas y los dos
cociclos conjuntos
\begin{equation}
\label{eq:continuo-comun-update-triadas-carry-conjunto}
 \begin{gathered}
 q_{K+1}=(u_{K+1}^\pi+u_{K+1}^e+u_{K+1}^\varphi)\bmod3,
 \quad
 c_{K+1}^{\rm joint}=\frac{u_{K+1}^\pi+u_{K+1}^e+u_{K+1}^\varphi-q_{K+1}}3,\\
 a_{K+1}=(u_{K+1}^\pi+u_{K+1}^e-u_{K+1}^\varphi)\bmod3,
 \quad
 h_{K+1}^{\rm joint}=\frac{u_{K+1}^\pi+u_{K+1}^e-u_{K+1}^\varphi-a_{K+1}}3,\\
 \boldsymbol\delta_{K+1}
 =\boldsymbol\delta_K\star
 \bigl((d_{K+1,\chi}^{\rm dec})_\chi,
       c_{K+1}^{\rm joint},h_{K+1}^{\rm joint}\bigr).
 \end{gathered}
\end{equation}
Así \(C_{K+1}\), \(\boldsymbol\delta_{K+1}\) y
\(\boldsymbol x_{K+1}\) tienen acciones materiales, no sólo nombres.
Las demás coordenadas se actualizan simultáneamente por
\begin{align}
 p_{K+1}&=p_K\star B_{K+1},
 \label{eq:continuo-comun-update-prefijo-tpk}\\
 \kappa(\gamma^{\rm TPK}_{0:K+1})
 &=\kappa(\varepsilon^{\rm TPK}_{K+1})
   +\mathsf C_{\varepsilon^{\rm TPK}_{K+1}}
    \kappa(\gamma^{\rm TPK}_{0:K}),
 \label{eq:continuo-comun-update-carry-lector-tpk}\\
 c_{K+1}^{\rm comp}
 &:=(\boldsymbol\delta_{K+1},\boldsymbol x_{K+1},
       \kappa(\gamma^{\rm TPK}_{0:K+1})),
 \label{eq:continuo-comun-update-carry-compacto-tpk}\\
 \Sigma_{K+1}
 &=\Sigma_B\bigl(
 B_{K+1},c_{K+1}^{\rm comp},
 \partial\gamma^{\rm TPK}_{0:K}
 +\partial\varepsilon^{\rm TPK}_{K+1},
 \omega_{K+1},\ell_{K+1}\bigr),
 \label{eq:continuo-comun-update-firma-tpk}\\
 W_{K+1}
 &=\bigl((u_{K+1}^\chi,u_{K+1}^\chi A_W)\bigr)_\chi,
 &g(K+1)&=g(K)+[1]_9.
 \label{eq:continuo-comun-update-witt-fase-tpk}
\end{align}
Como \(A_H\) es invertible módulo tres y \(A_W^2=-I_6\), las actualizaciones
residual y de sombra no colapsan bloques distintos. El truncamiento conserva
las secuencias completas
\[
 (B_0,\ldots,B_K),\ (C_0,\ldots,C_K),\
 (p_0,\ldots,p_K),\ (c_0,\ldots,c_K),\
 (\Sigma_0,\ldots,\Sigma_K),\ (W_0,\ldots,W_K)
\]
y elimina simultáneamente su última entrada. Por ello el lector
\(\mathfrak r_9\) conmuta con el borrado de un bloque.

\begin{proposition}[Alcance universal y alcance certificado del no reinicio]
\label{prop:continuo-comun-no-reinicio-alcance}
En toda vuelta horizontal de nueve pasos se conservan
\[
 g(K+9)=g(K),\qquad
 q_{\rm mem}(K+9)=q_{\rm mem}(K)+1,
\]
y el prefijo se prolonga con nueve emisiones, por lo que el estado no se
reinicia. En la sección conjunta certificada
\(\chi\in\{\pi,e,\varphi\}\), \(0\leq K<387\), se cumple además
\begin{equation}
\label{eq:continuo-comun-no-reinicio-certificado}
 (p_K,C_K)\neq(p_{K+9},C_{K+9}),\qquad
 (B_K,\mathbf N_K,W_K)\neq
 (B_{K+9},\mathbf N_{K+9},W_{K+9}),
 \qquad
 \mathbf N_K:=(N_K^\pi,N_K^e,N_K^\varphi).
\end{equation}
\end{proposition}

\begin{proof}
Las dos identidades universales son la iteración exacta de
\eqref{eq:continuo-comun-accion-fase-tpk}; el prefijo conserva las nueve
entradas añadidas por \eqref{eq:continuo-comun-update-prefijo-tpk}. Para la
sección conjunta, la recurrencia
\eqref{eq:continuo-comun-update-cilindro-indice} compara los \(387\) pares
nonádicos de cada uno de los tres canales y produce
\eqref{eq:continuo-comun-no-reinicio-certificado}; ésta es la afirmación
certificada en
\cref{thm:generacion-monodromica-conjunta-rev7}. Finalmente,
\(u\mapsto(u,uA_W)\) es inyectiva, de modo que el cambio de bloque conserva
también el cambio de su sombra completa. No se extiende esta última
desigualdad a una prolongación neutra arbitraria fuera de la sección
certificada.
\end{proof}

\begin{definition}[Estado APP--TRIT--TPK de profundidad \(k\)]
\label{def:continuo-comun-estado}
Para una ruta \(\gamma=\epsilon_1\cdots\epsilon_k\) que nace en una célula
APP \(a_0\), el estado enriquecido común es
\begin{equation}
\label{eq:continuo-comun-estado}
 \widehat x_k(\gamma):=\Bigl(
  a_0,\ldots,a_k;
  \operatorname{Sh}_0,\ldots,\operatorname{Sh}_k;
  \mathbf{res}_k,\mathbf{quo}_k;
  \mathbf{or}_k,\mathbf\kappa_k;
  \mathsf Q_k;g_k;q^{(9)}_k;\beta_k;m_k;\partial\gamma;\gamma;
  \operatorname{Led}(\gamma)
 \Bigr),
\end{equation}
donde
\begin{align*}
 \mathbf{res}_k
 &:=\bigl(
  (\operatorname{res}^{\Sigma}(\epsilon_j),
   \operatorname{res}^{\Pi}(\epsilon_j))
  \bigr)_{1\leq j\leq k},\\
 \mathbf{quo}_k
 &:=\bigl(
  (\operatorname{quo}^{\Sigma}(\epsilon_j),
   \operatorname{quo}^{\Pi}(\epsilon_j))
  \bigr)_{1\leq j\leq k},\\
 \mathbf{or}_k
 &:=\bigl(\operatorname{or}(\epsilon_j)\bigr)_{1\leq j\leq k},
 \qquad
 \mathbf\kappa_k
 :=\bigl(\kappa(\epsilon_j)\bigr)_{1\leq j\leq k},\\
 m_k&:=m(\gamma;m_0).
\end{align*}
Las hojas intermedias se obtienen por
\(\operatorname{Sh}_{j}=U_{\epsilon_j}\operatorname{Sh}_{j-1}\).
\end{definition}

La presencia simultánea de la ruta y del ledger impide identificar palabras
distintas que terminan en la misma célula.  El estado no es una lista de
cardinales: cada coordenada tiene las ecuaciones de tipo y composición de las
secciones anteriores.

\begin{definition}[Dominio total de profundidad finita]
\label{def:continuo-comun-dominio-total}
Sea
\begin{equation}
\label{eq:continuo-comun-dominio-total}
 \widehat{\mathcal X}^{\rm tot}_k
 :=\left\{
  \widehat x_k(\epsilon_1\cdots\epsilon_k):
  \begin{array}{l}
   a_0\in\mathsf A_0,\\
   \epsilon_j\in\mathsf E_{j-1}(a_{j-1}),\\
   s(\epsilon_j)=a_{j-1},\ t(\epsilon_j)=a_j
  \end{array}
 \right\}.
\end{equation}
No se añade a \eqref{eq:continuo-comun-dominio-total} ninguna condición de
una publicación posterior.  Ser un elemento del dominio significa
exactamente haber sido producido por las acciones totales
\eqref{eq:continuo-comun-residuo-cociente-app},
\eqref{eq:continuo-comun-trit-euclid} y
\eqref{eq:continuo-comun-accion-fase-tpk}--
\eqref{eq:continuo-comun-accion-memoria}.
\end{definition}

El truncamiento común borra la última emisión y todas y sólo las entradas del
ledger que ésta produjo:
\begin{equation}
\label{eq:continuo-comun-truncamiento-total}
 \rho_{k+1,k}:\widehat{\mathcal X}^{\rm tot}_{k+1}
 \longrightarrow\widehat{\mathcal X}^{\rm tot}_k,
 \qquad
 \rho_{k+1,k}\widehat x_{k+1}
 =\widehat x_k.
\end{equation}
Sea \(\varepsilon^0_k(x)\) la emisión neutra
\eqref{eq:continuo-comun-emision-neutra-app} aplicada al extremo APP de
\(x\).  La prolongación neutra define
\begin{equation}
\label{eq:continuo-comun-seccion-total}
 \eta_k:\widehat{\mathcal X}^{\rm tot}_k
 \longrightarrow\widehat{\mathcal X}^{\rm tot}_{k+1},
 \qquad
 \eta_k(x):=(\varepsilon^0_k(x))_*x,
 \qquad
 \rho_{k+1,k}\eta_k=I.
\end{equation}

\begin{proposition}[Finitud, no vacuidad y sobreyectividad]
\label{prop:continuo-comun-total-finito}
Para todo \(k\), el conjunto
\(\widehat{\mathcal X}^{\rm tot}_k\) es finito y no vacío, y
\(\rho_{k+1,k}\) es sobreyectivo.
\end{proposition}

\begin{proof}
APP produce un conjunto finito no vacío en profundidad cero.  En cada paso,
cada estado tiene un conjunto finito de emisiones por
\eqref{eq:continuo-comun-out-no-vacio-app}; por inducción, sólo hay un número
finito de palabras de longitud \(k\).  La sucesión de prolongaciones neutras
produce al menos una palabra de cada longitud.  Finalmente,
\eqref{eq:continuo-comun-seccion-total} es una inversa por la derecha del
truncamiento, de modo que éste es sobreyectivo.
\end{proof}

\section{Grupoide finito de simetrías reversibles}
\label{sec:continuo-comun-groupoide}

Las operaciones reversibles de APP, TRIT y TPK actúan sobre el estado entero.
Sea \(\mathsf U\) el grupo de familias coherentes
\(g=(g_k)_{k\geq0}\) generadas por esas operaciones, donde
\begin{equation}
\label{eq:continuo-comun-simetria-coherente}
 \rho_{k+1,k}g_{k+1}=g_k\rho_{k+1,k},
 \qquad
 g_{k+1}\eta_k=\eta_kg_k.
\end{equation}
Denotamos por \(\mathsf U_k\) su imagen en
\(\operatorname{Sym}(\widehat{\mathcal X}^{\rm tot}_k)\).  Esta imagen es
finita porque actúa sobre un conjunto finito.  La acción de \(g_k\) sobre
\begin{equation*}
 \widehat x=(a_{\leq k};\operatorname{Sh}_{\leq k};
 \mathbf{res},\mathbf{quo};\mathbf{or},\mathbf\kappa;
 \mathsf Q_k;g_k;q_k^{(9)};\beta_k;m;\partial\gamma;\gamma;
 \operatorname{Led}(\gamma))
\end{equation*}
es
\begin{equation}
\label{eq:continuo-comun-accion-reversible}
\begin{aligned}
 g_k\cdot\widehat x:=\bigl(&
 g_k a_{\leq k};
 U_{g_k}\operatorname{Sh}_{\leq k};
 \mathbf{res}(g_k\gamma),\mathbf{quo}(g_k\gamma);
 \mathbf{or}(g_k\gamma),\mathbf\kappa(g_k\gamma);\\
 &\mathsf Q(g_k\gamma);g(g_k\gamma);q^{(9)}(g_k\gamma);
 \beta(g_k\gamma);
 \mathsf U_{g_k}(m);
 (g_k)_*\partial\gamma;
 g_k\gamma;
 \operatorname{Led}(g_k\gamma)\bigr).
\end{aligned}
\end{equation}
Todos los términos del lado derecho se recalculan mediante las reglas APP,
TRIT y TPK; no se copian como etiquetas independientes.

\begin{definition}[Grupoide de acción común]
\label{def:continuo-comun-groupoide}
El grupoide finito de profundidad \(k\) es
\begin{equation}
\label{eq:continuo-comun-groupoide}
 \mathfrak G_k
 :=\mathsf U_k\ltimes\widehat{\mathcal X}^{\rm tot}_k.
\end{equation}
Sus objetos son los estados comunes.  Una flecha
\((g_k,\widehat x)\) tiene fuente \(\widehat x\), destino
\(g_k\cdot\widehat x\), identidad \((I,\widehat x)\) e inversa
\((g_k^{-1},g_k\cdot\widehat x)\).
\end{definition}

\begin{lemma}[La acción conserva todas las ecuaciones de estado]
\label{lem:continuo-comun-groupoide-conserva}
La acción \eqref{eq:continuo-comun-accion-reversible} conserva la pareja de
hojas, las divisiones APP, la orientación TRIT, el carry, el odómetro de
fase, la memoria vertical, su reducción dodecafásica, la memoria afín, la
frontera, la ruta y el ledger, en sus tipos respectivos.  Además,
\begin{equation}
\label{eq:continuo-comun-groupoide-accion}
 I\cdot\widehat x=\widehat x,
 \qquad
 (g_{2,k}g_{1,k})\cdot\widehat x
 =g_{2,k}\cdot(g_{1,k}\cdot\widehat x).
\end{equation}
\end{lemma}

\begin{proof}
Las divisiones APP son únicas, por lo que la acción sobre una emisión
determina de nuevo sus residuos y cocientes.  La descomposición equilibrada
\eqref{eq:continuo-comun-trit-euclid} es también única.  La composición del
carry es asociativa por
\cref{lem:continuo-comun-asociatividad-carry}; la acción sobre memoria lo es
por \cref{lem:continuo-comun-accion-memoria}; y la frontera es functorial por
\eqref{eq:continuo-comun-frontera-telescopica}.  Finalmente, la acción sobre
la palabra ordenada induce la acción sobre su ledger entrada por entrada.
La coherencia con el odómetro recalcula \(Q\), \(g\) y \(q^{(9)}\) a partir
de la ruta transformada y conserva
\eqref{eq:continuo-comun-accion-fase-tpk}. Estas identidades prueban
\eqref{eq:continuo-comun-groupoide-accion}.
\end{proof}

\section{Árbol finito de todas las prolongaciones}
\label{sec:continuo-comun-arbol}

Para \(x\in\widehat{\mathcal X}^{\rm tot}_k\), sea
\(\operatorname{Out}(x)\) el conjunto de todas las emisiones TPK admisibles
de longitud uno desde el extremo de \(x\).  Una emisión pertenece a
\(\operatorname{Out}(x)\) sólo cuando su acción actualiza conjuntamente todas
las coordenadas de \eqref{eq:continuo-comun-estado}.  En particular,
\begin{equation}
\label{eq:continuo-comun-out-finito}
 1\leq d(x):=\#\operatorname{Out}(x)<\infty,
 \qquad
 \varepsilon^0_k(x)\in\operatorname{Out}(x).
\end{equation}

\begin{definition}[Árbol truncado de prolongaciones]
\label{def:continuo-comun-arbol-finito}
Dados \(N\geq k\) y
\(x\in\widehat{\mathcal X}^{\rm tot}_k\), definimos
\(\mathscr T_{k,N}(x)\) como el árbol enraizado cuyos vértices de altura
\(j\), \(0\leq j\leq N-k\), son las palabras
\begin{equation}
\label{eq:continuo-comun-vertice-arbol}
 h=(\epsilon_{k+1},\ldots,\epsilon_{k+j}),
 \qquad
 \epsilon_{n+1}\in\operatorname{Out}(x_n),
 \qquad
 x_{n+1}=(\epsilon_{n+1})_*x_n,
\end{equation}
con raíz la palabra vacía.  Una arista añade una emisión al final.  Dos
palabras que alcanzan el mismo estado terminal siguen siendo vértices
distintos si sus ledgers difieren.
\end{definition}

La última condición es esencial: el árbol registra genealogías y no el
cociente que sólo recuerda el destino.  Su nivel terminal es
\begin{equation}
\label{eq:continuo-comun-terminal-finito}
 \operatorname{Term}_{k,N}(x)
 :=\{h\in\mathscr T_{k,N}(x):|h|=N-k\}.
\end{equation}
Cada \(h\) conserva el estado terminal completo
\begin{equation}
\label{eq:continuo-comun-registro-terminal}
 \operatorname{ter}_{N}(h)
 :=\bigl(x_N,m_N,\partial(\gamma h),
          \gamma h,\operatorname{Led}(\gamma h)\bigr).
\end{equation}

\begin{proposition}[Finitud y no vacuidad a todo horizonte]
\label{prop:continuo-comun-arbol-finito-no-vacio}
Para todo \(x\), \(k\) y \(N\geq k\), el árbol
\(\mathscr T_{k,N}(x)\) es finito y
\(\operatorname{Term}_{k,N}(x)\neq\varnothing\).
\end{proposition}

\begin{proof}
La raíz tiene un número finito y positivo de sucesores por
\eqref{eq:continuo-comun-out-finito}.  Aplicando el mismo argumento en cada
nivel se obtiene, por inducción, finitud a todo horizonte finito.  La palabra
formada por emisiones neutras
\((\varepsilon^0_k,\varepsilon^0_{k+1},\ldots,
\varepsilon^0_{N-1})\) es un vértice terminal, por lo que
el nivel \eqref{eq:continuo-comun-terminal-finito} no es vacío.
\end{proof}

El borrado de la última emisión define
\begin{equation}
\label{eq:continuo-comun-proyeccion-terminal}
 \pi_{N+1,N}:\operatorname{Term}_{k,N+1}(x)
 \longrightarrow\operatorname{Term}_{k,N}(x).
\end{equation}

\begin{lemma}[Sobreyectividad terminal]
\label{lem:continuo-comun-terminal-sobreyectivo}
La aplicación \(\pi_{N+1,N}\) es sobreyectiva para todo \(N\geq k\).
\end{lemma}

\begin{proof}
Sea \(h\in\operatorname{Term}_{k,N}(x)\), de estado terminal \(v_h\).
Añadir la emisión \(\varepsilon_N^0(v_h)\) produce
\(h'\in\operatorname{Term}_{k,N+1}(x)\) y
\(\pi_{N+1,N}(h')=h\).
\end{proof}

\section{Conexión nonádica: fase, holonomía y memoria}
\label{sec:continuo-comun-gamma9-total}

Una palabra de nueve emisiones no se denomina holonomía sólo por tener
longitud nueve. Primero se conserva la conexión horizontal determinada por
el funtor de fase. Para \(r\in C_9\), definimos
\begin{equation}
\label{eq:continuo-comun-relacion-horizontal-fase}
 \begin{split}
 \mathscr R_{r,k}:=\bigl\{
 (x,\epsilon,\epsilon_*x):\;&
 x\in\widehat{\mathcal X}^{\rm tot}_k,\quad
 \epsilon\in\operatorname{Out}(x),\\
 &\mathsf Q_k(x)=(w,r),\quad
 \mathsf Q_{k+1}(\epsilon_*x)=\sigma_9(w,r)
 \bigr\}.
 \end{split}
\end{equation}
La última igualdad no recorta el dominio: es la acción total
\eqref{eq:continuo-comun-accion-fase-tpk} de toda emisión TPK elemental.
La conexión transporta simultáneamente hoja, memoria y ledger:
\begin{equation}
\label{eq:continuo-comun-conexion-tpk-generador}
 \nabla_{\rm TPK}(\epsilon):=
 \bigl(\operatorname{ph}(\epsilon),U_\epsilon,
       \mathsf U_\epsilon,\lambda(\epsilon)\bigr).
\end{equation}
Su composición se calcula mediante
\eqref{eq:continuo-comun-transporte-hojas-funtorial},
\eqref{eq:continuo-comun-accion-memoria-composicion} y la concatenación del
ledger.

El ápice de la vuelta nonádica es ahora el conjunto de lifts horizontales
fase-compatibles
\begin{equation}
\label{eq:continuo-comun-gamma9-apice}
 \begin{split}
 \mathsf Z^{\Gamma}_{9,k}:=\coprod_{r\in C_9}\bigl\{
 (x,\epsilon_1,\ldots,\epsilon_9):\;&x_0=x,\quad
 x_i=(\epsilon_i)_*x_{i-1},\\
 &(x_{i-1},\epsilon_i,x_i)
 \in\mathscr R_{r+i-1,k+i-1},\quad 1\leq i\leq9
 \bigr\},
 \end{split}
\end{equation}
 donde los subíndices de \(r+i-1\) se leen en \(C_9\). El sumando está
determinado de forma única por \(r=g_k(x)\), de modo que el coproducto no
duplica testigos. Los mapas fuente y
destino son
\begin{equation}
\label{eq:continuo-comun-gamma9-span}
 \widehat{\mathcal X}^{\rm tot}_k
 \xleftarrow{\ s_{9,k}\ }
 \mathsf Z^{\Gamma}_{9,k}
 \xrightarrow{\ t_{9,k}\ }
 \widehat{\mathcal X}^{\rm tot}_{k+9},
 \qquad
 s_{9,k}(x,\boldsymbol\epsilon)=x,\quad
 t_{9,k}(x,\boldsymbol\epsilon)
 =(\epsilon_9\cdots\epsilon_1)_*x.
\end{equation}
La holonomía relacional es
\begin{equation}
\label{eq:continuo-comun-gamma9-relacion}
 \Gamma_{9,k}:=
 \{(s_{9,k}z,t_{9,k}z):z\in\mathsf Z^{\Gamma}_{9,k}\}
 =\operatorname{Hol}_{C_9}(\nabla_{\rm TPK}).
\end{equation}

\begin{lemma}[El nivel nueve del árbol realiza la holonomía]
\label{lem:continuo-comun-gamma9-arbol}
La aplicación que olvida los estados intermedios ya determinados por las
emisiones induce una biyección
\begin{equation}
\label{eq:continuo-comun-gamma9-apice-arbol}
 \mathsf Z^{\Gamma}_{9,k}
 \xrightarrow{\ \cong\ }
 \coprod_{x\in\widehat{\mathcal X}^{\rm tot}_k}
 \operatorname{Term}_{k,k+9}(x).
\end{equation}
\end{lemma}

\begin{proof}
Cada arista de \(\operatorname{Out}(x_i)\) satisface
\eqref{eq:continuo-comun-accion-fase-tpk}; por ello una palabra terminal de
longitud nueve determina los nueve elementos de
\eqref{eq:continuo-comun-relacion-horizontal-fase}. Recíprocamente, un lift
horizontal de \eqref{eq:continuo-comun-gamma9-apice} es una palabra
componible del árbol. Las dos operaciones conservan las emisiones y son
inversas. Así \eqref{eq:continuo-comun-gamma9-apice-arbol} es un teorema
derivado de la conexión, no la definición de ésta.
\end{proof}

\begin{proposition}[Retorno de fase, avance de memoria y no reinicio]
\label{prop:continuo-comun-gamma9-retorno-memoria}
Para todo \(z=(x,\epsilon_1,\ldots,\epsilon_9)\) del ápice,
\begin{align}
 g(t_{9,k}z)&=g(s_{9,k}z),
 \label{eq:continuo-comun-gamma9-retorno-fase}\\
 q^{(9)}(t_{9,k}z)&=q^{(9)}(s_{9,k}z)+1,
 \label{eq:continuo-comun-gamma9-avance-memoria}\\
 \operatorname{Led}(t_{9,k}z)
 &=\operatorname{Led}(s_{9,k}z)\star
   \bigl(\lambda(\epsilon_1),\ldots,\lambda(\epsilon_9)\bigr).
 \label{eq:continuo-comun-gamma9-ledger}
\end{align}
En particular, la identificación de fase
\(g(t_{9,k}z)=g(s_{9,k}z)\) no induce una identificación de estados: tras
transportar ambos registros al mismo tipo de comparación, sus memorias
verticales difieren en una unidad. Retorna la fase publicada, no el estado
enriquecido.
\end{proposition}

\begin{proof}
Las dos primeras identidades son
\eqref{eq:continuo-comun-retorno-fase-avance-memoria}. La tercera es la
definición de concatenación del ledger. Como la memoria vertical difiere en
una unidad, los estados completos son distintos aunque alguna coordenada
visible adicional pudiera coincidir.
\end{proof}

\begin{proposition}[Compresión isométrica de la misma holonomía]
\label{prop:continuo-comun-gamma9-compresion-isometrica}
Sea \(U_{\Gamma_9}:\mathcal H_{\rm full}\to\mathcal H_{\rm full}\) una
realización isométrica de la holonomía anterior y sea
\(J:\mathcal H_{\rm vis}\to\mathcal H_{\rm full}\) la inclusión isométrica
del registro visible, con \(E:=J^*\) y \(P:=JE\). Definimos
\begin{equation}
\label{eq:continuo-comun-gamma9-compresion-datos}
 T:=EU_{\Gamma_9}J,\qquad
 \eta:=(I-P)U_{\Gamma_9}J.
\end{equation}
Entonces
\begin{equation}
\label{eq:continuo-comun-gamma9-compresion}
 \boxed{\;
 U_{\Gamma_9}J=JT+\eta,\qquad
 I-T^*T=\eta^*\eta.
 \;}
\end{equation}
\end{proposition}

\begin{proof}
Insertando \(I=P+(I-P)\) delante de \(U_{\Gamma_9}J\) se obtiene
\[
 U_{\Gamma_9}J
 =JEU_{\Gamma_9}J+(I-P)U_{\Gamma_9}J
 =JT+\eta.
\]
Los dos sumandos pertenecen a subespacios ortogonales. Como
\(U_{\Gamma_9}\) y \(J\) son isometrías,
\[
 I=(U_{\Gamma_9}J)^*(U_{\Gamma_9}J)
   =T^*T+\eta^*\eta,
\]
que es la segunda identidad.
\end{proof}

\begin{proposition}[Totalidad y equivariancia de la holonomía]
\label{prop:continuo-comun-gamma9-total-natural}
El mapa \(s_{9,k}\) es sobreyectivo. Las nueve prolongaciones neutras forman
la sección
\begin{equation}
\label{eq:continuo-comun-gamma9-seccion-neutra}
 \begin{gathered}
 x_0:=x,\qquad
 \varepsilon^0_{k+i}:=\varepsilon^0_{k+i}(x_i),\qquad
 x_{i+1}:=(\varepsilon^0_{k+i})_*x_i\quad(0\leq i<9),\\
 \eta^{(9)}_k(x):=
 (x,\varepsilon^0_k,\ldots,\varepsilon^0_{k+8})
 \in\mathsf Z^{\Gamma}_{9,k},
 \qquad s_{9,k}\eta^{(9)}_k=I.
 \end{gathered}
\end{equation}
Toda simetría coherente \(a\) cuya acción en el odómetro satisface
\(\mathsf Q\,a=\bar a\,\mathsf Q\) y
\(\bar a\,\sigma_9=\sigma_9\,\bar a\) induce
\[
 a^Z(x,\epsilon_1,\ldots,\epsilon_9)
 :=(a x,a\epsilon_1,\ldots,a\epsilon_9)
\]
y verifica
\begin{equation}
\label{eq:continuo-comun-gamma9-naturalidad}
 s_{9,k}a^Z=a_ks_{9,k},
 \qquad
 t_{9,k}a^Z=a_{k+9}t_{9,k}.
\end{equation}
\end{proposition}

\begin{proof}
La palabra neutra existe en toda fuente y satisface la acción de fase
\eqref{eq:continuo-comun-accion-fase-tpk}; esto prueba la sección y la
sobreyectividad. La condición
\(\bar a\sigma_9=\sigma_9\bar a\) envía relaciones horizontales a relaciones
horizontales. Actuar entrada por entrada conserva composición, carry,
memoria, frontera y ledger; leer el primer o el último estado antes o después
de actuar da \eqref{eq:continuo-comun-gamma9-naturalidad}.
\end{proof}

\section{Supervivencia total y multisección completa}
\label{sec:continuo-comun-supervivencia}

\begin{definition}[Supervivencia producida]
\label{def:continuo-comun-supervivencia}
Un estado \(x\in\widehat{\mathcal X}^{\rm tot}_k\) sobrevive cuando
\begin{equation}
\label{eq:continuo-comun-supervivencia}
 \operatorname{Term}_{k,N}(x)\neq\varnothing
 \qquad\text{para todo }N\geq k.
\end{equation}
Denotamos por \(\operatorname{Surv}_k\) el conjunto de tales estados.
\end{definition}

\begin{theorem}[El dominio total es superviviente]
\label{thm:continuo-comun-total-superviviente}
Para toda profundidad,
\begin{equation}
\label{eq:continuo-comun-surv-igual-total}
 \operatorname{Surv}_k=\widehat{\mathcal X}^{\rm tot}_k.
\end{equation}
En particular, la supervivencia no define un subconjunto supuesto antes de
construir el estado.
\end{theorem}

\begin{proof}
La inclusión \(\operatorname{Surv}_k\subseteq
\widehat{\mathcal X}^{\rm tot}_k\) pertenece a la definición.  La inclusión
opuesta es exactamente la no vacuidad a todo horizonte probada en
\cref{prop:continuo-comun-arbol-finito-no-vacio}.
\end{proof}

\begin{definition}[Multisección de prolongaciones completas]
\label{def:continuo-comun-multiseccion}
La multisección completa sobre \(x\) es el límite inverso
\begin{equation}
\label{eq:continuo-comun-multiseccion}
 \operatorname{Msect}(x)
 :=\varprojlim_{N\geq k}
 \bigl(\operatorname{Term}_{k,N}(x),\pi_{N+1,N}\bigr).
\end{equation}
Sus elementos son todas las historias infinitas compatibles que prolongan
\(x\).  Ninguna de ellas recibe prioridad dentro de
\eqref{eq:continuo-comun-multiseccion}.
\end{definition}

\begin{theorem}[No vacuidad material de la multisección]
\label{thm:continuo-comun-multiseccion-no-vacia}
Para todo \(x\in\widehat{\mathcal X}^{\rm tot}_k\),
\begin{equation}
\label{eq:continuo-comun-multiseccion-no-vacia}
 \operatorname{Msect}(x)\neq\varnothing.
\end{equation}
Además, cada coordenada finita de una historia de
\(\operatorname{Msect}(x)\) conserva las dos hojas, orientación, carry,
memoria, frontera, ruta y ledger del estado que prolonga.
\end{theorem}

\begin{proof}
Los conjuntos terminales son finitos y no vacíos por
\cref{prop:continuo-comun-arbol-finito-no-vacio}, y sus aplicaciones de
enlace son sobreyectivas por
\cref{lem:continuo-comun-terminal-sobreyectivo}.  Se puede construir un
elemento compatible recursivamente: se comienza en la raíz y, en cada paso,
se conserva una prolongación del prefijo anterior.  La prolongación neutra
garantiza que el proceso nunca se detenga.  La compatibilidad de las
coordenadas se sigue de que cada arista del árbol actúa mediante la única
acción de estado definida en
\eqref{eq:continuo-comun-estado}.
\end{proof}

\section{Medida canónica generada por el árbol}
\label{sec:continuo-comun-medida}

La medida no se introduce como un peso exterior sobre las salidas.  Se obtiene
contando las emisiones admisibles en cada vértice antes de elegir una
continuación.

\begin{definition}[Peso local y medida de horizonte]
\label{def:continuo-comun-medida-finita}
Para \(v\in\widehat{\mathcal X}^{\rm tot}_j\) y
\(\epsilon\in\operatorname{Out}(v)\), ponemos
\begin{equation}
\label{eq:continuo-comun-peso-local}
 p_v(\epsilon):=\frac{1}{d(v)}\in\mathbb Q_{>0}.
\end{equation}
Si
\(h=(\epsilon_{k+1},\ldots,\epsilon_N)\in
\operatorname{Term}_{k,N}(x)\) y \(x_j\) es el prefijo que precede a
\(\epsilon_{j+1}\), definimos
\begin{equation}
\label{eq:continuo-comun-medida-horizonte}
 \mu_{x,N}(h)
 :=\prod_{j=k}^{N-1}p_{x_j}(\epsilon_{j+1})
 =\prod_{j=k}^{N-1}\frac{1}{d(x_j)}.
\end{equation}
\end{definition}

\begin{lemma}[Normalización]
\label{lem:continuo-comun-medida-normalizada}
Para todo horizonte \(N\geq k\),
\begin{equation}
\label{eq:continuo-comun-medida-suma-uno}
 \sum_{h\in\operatorname{Term}_{k,N}(x)}\mu_{x,N}(h)=1.
\end{equation}
\end{lemma}

\begin{proof}
En altura uno,
\(\sum_{\epsilon\in\operatorname{Out}(x)}d(x)^{-1}=1\).  Supuesta la
normalización a altura \(N-k\), cada prefijo terminal \(h\), de estado final
\(v_h\), aporta en el nivel siguiente
\[
 \sum_{\epsilon\in\operatorname{Out}(v_h)}
 \mu_{x,N}(h)p_{v_h}(\epsilon)
 =\mu_{x,N}(h).
\]
Sumar sobre todos los prefijos prueba la identidad por inducción.
\end{proof}

\begin{proposition}[Consistencia proyectiva]
\label{prop:continuo-comun-medida-consistente}
Las medidas finitas satisfacen
\begin{equation}
\label{eq:continuo-comun-medida-pushforward}
 (\pi_{N+1,N})_*\mu_{x,N+1}=\mu_{x,N}.
\end{equation}
\end{proposition}

\begin{proof}
Para \(h\in\operatorname{Term}_{k,N}(x)\), la masa de su fibra es
\begin{align*}
 \sum_{\pi_{N+1,N}(h')=h}\mu_{x,N+1}(h')
 &=\mu_{x,N}(h)
   \sum_{\epsilon\in\operatorname{Out}(v_h)}p_{v_h}(\epsilon)\\
 &=\mu_{x,N}(h).
\end{align*}
Ésta es precisamente la igualdad de medidas puntuales que equivale a
\eqref{eq:continuo-comun-medida-pushforward}.
\end{proof}

Los cilindros de la multisección son
\begin{equation}
\label{eq:continuo-comun-cilindro-msect}
 [h]_{x,N}
 :=\{\omega\in\operatorname{Msect}(x):
       \omega|_N=h\}.
\end{equation}

\begin{theorem}[Medida canónica de prolongaciones]
\label{thm:continuo-comun-medida-canonica}
Existe una única medida de cilindros \(\mu_x\) sobre
\(\operatorname{Msect}(x)\) tal que
\begin{equation}
\label{eq:continuo-comun-medida-cilindro}
 \mu_x([h]_{x,N})=\mu_{x,N}(h).
\end{equation}
Está normalizada, tiene soporte en toda la multisección y se obtiene sólo de
los conjuntos finitos de prolongaciones APP--TRIT--TPK.
\end{theorem}

\begin{proof}
La consistencia de
\cref{prop:continuo-comun-medida-consistente} hace que
\eqref{eq:continuo-comun-medida-cilindro} sea independiente del horizonte en
el que se represente un cilindro.  La normalización procede de
\cref{lem:continuo-comun-medida-normalizada}.  Dos cilindros o bien son
disjuntos o bien uno contiene al otro; por tanto, las masas finitas definen
de manera única una medida aditiva en el álgebra de cilindros.  La
multisección es un límite de conjuntos finitos discretos y los cilindros
forman su base compacta; la medida se prolonga de manera única a la
\(\sigma\)-álgebra generada por ellos.  Finalmente,
\(p_v(\epsilon)>0\) para toda emisión admisible, de modo que todo cilindro no
vacío tiene masa positiva y el soporte es la multisección completa.
\end{proof}

\section{Naturalidad bajo simetrías y truncamientos}
\label{sec:continuo-comun-naturalidad}

Una familia coherente \(g=(g_j)_{j\geq0}\in\mathsf U\) transporta una
emisión saliente de profundidad \(j\) mediante
\begin{equation}
\label{eq:continuo-comun-accion-out}
 (g_{j+1},g_j)_*:\operatorname{Out}(x)\longrightarrow
      \operatorname{Out}(g_j\cdot x),
 \qquad
 \epsilon\longmapsto g_{j+1}\epsilon g_j^{-1}.
\end{equation}
Las relaciones TPK hacen de \eqref{eq:continuo-comun-accion-out} una biyección
que conserva multiplicidades.  Actuando en cada entrada se obtiene un
isomorfismo de árboles
\begin{equation}
\label{eq:continuo-comun-isomorfismo-arbol}
 g_{*,N}:\mathscr T_{k,N}(x)
 \xrightarrow{\ \cong\ }
 \mathscr T_{k,N}(g_k\cdot x).
\end{equation}

\begin{theorem}[Naturalidad del árbol, la multisección y la medida]
\label{thm:continuo-comun-naturalidad}
Para toda simetría reversible coherente \(g=(g_j)_{j\geq0}\),
\begin{align}
 \pi_{N+1,N}g_{*,N+1}
 &=g_{*,N}\pi_{N+1,N},
 \label{eq:continuo-comun-naturalidad-terminal}\\
 g_*\operatorname{Msect}(x)
 &=\operatorname{Msect}(g_k\cdot x),
 \label{eq:continuo-comun-naturalidad-msect}\\
 (g_{*,N})_*\mu_{x,N}
 &=\mu_{g_k\cdot x,N},
 \qquad
 (g_*)_*\mu_x=\mu_{g_k\cdot x}.
 \label{eq:continuo-comun-naturalidad-medida}
\end{align}
\end{theorem}

\begin{proof}
La primera igualdad expresa que actuar entrada por entrada y borrar la última
entrada conmutan.  Al pasar al límite inverso se obtiene
\eqref{eq:continuo-comun-naturalidad-msect}.  Como
\eqref{eq:continuo-comun-accion-out} es biyectiva,
\(d(g_j\cdot v)=d(v)\); por
\eqref{eq:continuo-comun-medida-horizonte}, una palabra y su imagen tienen el
mismo peso.  Esto prueba la primera igualdad de medidas de
\eqref{eq:continuo-comun-naturalidad-medida}.  La segunda se deduce en los
cilindros mediante \eqref{eq:continuo-comun-medida-cilindro}, y los cilindros
determinan la medida.
\end{proof}

Los truncamientos del estado y los del árbol forman asimismo el cuadrado
\begin{equation}
\label{eq:continuo-comun-cuadrado-truncamiento}
\begin{CD}
 \operatorname{Term}_{k,N+1}(x)
 @>{\pi_{N+1,N}}>>
 \operatorname{Term}_{k,N}(x)\\
 @V{\operatorname{ter}_{N+1}}VV
 @VV{\operatorname{ter}_{N}}V\\
 \widehat{\mathcal X}^{\rm tot}_{N+1}
 @>{\rho_{N+1,N}}>>
 \widehat{\mathcal X}^{\rm tot}_{N}.
\end{CD}
\end{equation}
La conmutatividad es literal: ambos recorridos borran la última emisión y su
última entrada de ledger.

\section{Límite inverso \(\widehat\Omega\) de historias completas y estado enriquecido común APP–TRIT–TPK}
\label{sec:continuo-comun-limite}

Las sobreyecciones \eqref{eq:continuo-comun-truncamiento-total} producen el
objeto de historias completas
\begin{equation}
\label{eq:continuo-comun-omega}
 \widehat\Omega
 :=\varprojlim_k
 \bigl(\widehat{\mathcal X}^{\rm tot}_k,\rho_{k+1,k}\bigr).
\end{equation}
Un elemento \(\widehat x=(\widehat x_k)_{k\geq0}\) contiene, en cada
profundidad, la misma pareja de hojas transportada, las orientaciones y
acarreos producidos, la memoria acumulada, las fronteras, la ruta y el ledger
completo.  Su fibra de continuaciones desde cualquier prefijo es
\(\operatorname{Msect}(\widehat x_k)\), equipada con
\(\mu_{\widehat x_k}\).

\begin{theorem}[Estado enriquecido común APP--TRIT--TPK]
\label{thm:continuo-comun-propietario}
El sistema
\begin{equation}
\label{eq:continuo-comun-propietario}
 \mathfrak C^{\rm gen}:=\Biggl(
 \begin{aligned}
 &\{\widehat{\mathcal X}^{\rm tot}_k,\rho_{k+1,k},\eta_k\}_{k\geq0},
   \widehat\Omega,\{\mathfrak G_k\}_{k\geq0},
   \{\mathscr T_{k,N}\}_{N\geq k},\\
 &\operatorname{Msect},\operatorname{Surv},\{\mu_x\},
   \{\mathfrak r_9(v_K)\}_{K\geq0},
   \mathcal K_{\rm ph},\widehat{\mathcal K}_{\rm ph},\\
 &\operatorname{Sh},\operatorname{or},\kappa,
   \mathsf Q,g,q^{(9)},\beta,m,\partial,\gamma,\operatorname{Led},\\
 &\{\mathsf Z^\Gamma_{9,k},s_{9,k},t_{9,k},\Gamma_{9,k}\}_{k\geq0}
 \end{aligned}
 \Biggr)
\end{equation}
está definido para todo estado producido por APP--TRIT--TPK\@.  Es no vacío,
natural bajo las simetrías reversibles y compatible con refinamiento.  En
particular:
\begin{enumerate}
 \item las dos hojas se transportan juntas;
 \item orientación y carry se calculan por la división TRIT;
 \item la memoria obedece la acción afín TPK;
 \item la fase obedece el odómetro TPK, vuelve tras nueve emisiones y la
       memoria vertical avanza una unidad;
 \item la realización Markov de esa vuelta satisface
       \(\mathcal K_{\rm ph}^9=E_+E_\times\) y renueva uniformemente la
       fibra completa de \(468=18\cdot26\) emisiones;
 \item el lector compacto conserva bloque, cilindro, prefijo, carry, firma,
       sombra de Witt y fase, con sus actualizaciones y truncamientos;
 \item frontera, ruta y ledger se componen sin borrar los extremos ni el
       orden de las emisiones;
 \item todo estado sobrevive y posee una multisección no vacía de
       prolongaciones completas;
 \item la medida canónica es proyectiva y tiene soporte en todas esas
       prolongaciones.
\end{enumerate}
\end{theorem}

\begin{proof}
La totalidad y la finitud proceden de
\cref{prop:continuo-comun-total-finito}.  Las identidades de hojas, carry y
memoria están probadas en
\cref{lem:continuo-comun-asociatividad-carry,lem:continuo-comun-accion-memoria}; la fase, la memoria vertical y la
holonomía nonádica, incluida su realización de fibra y su lector de estado,
están probadas en
\cref{prop:continuo-comun-gamma9-retorno-memoria,prop:continuo-comun-gamma9-total-natural,prop:continuo-comun-kph-heredada,prop:continuo-comun-gamma9-compresion-isometrica,prop:continuo-comun-no-reinicio-alcance}; la frontera conserva sus extremos por
\eqref{eq:continuo-comun-frontera-telescopica}.  La supervivencia total y la
no vacuidad de las multisecciones son
\cref{thm:continuo-comun-total-superviviente,thm:continuo-comun-multiseccion-no-vacia}.  La existencia, consistencia y
soporte de la medida son
\cref{prop:continuo-comun-medida-consistente,thm:continuo-comun-medida-canonica}.  Finalmente, la compatibilidad con las
simetrías y con los truncamientos se demuestra en
\cref{thm:continuo-comun-naturalidad} y
\eqref{eq:continuo-comun-cuadrado-truncamiento}.
\end{proof}

\section{Necesidad estructural y procedencia}
\label{sec:continuo-comun-falsadores}

\begin{proposition}[Falsadores del estado enriquecido común]
\label{prop:continuo-comun-falsadores}
Cualquiera de las operaciones siguientes sustituye
\(\mathfrak C^{\rm gen}\) por otro objeto y bloquea la transferencia de los
teoremas anteriores:
\begin{enumerate}
 \item elegir \(\Sigma\) o \(\Pi\) y descartar la otra hoja;
 \item conservar el residuo y borrar el cociente;
 \item conservar el signo TRIT y borrar su carry;
 \item reemplazar el ledger ordenado por su valor terminal;
 \item identificar dos palabras porque alcanzan la misma célula;
 \item escoger una continuación y sustituir por ella la multisección;
 \item declarar supervivencia sin exhibir prolongaciones a cada horizonte;
 \item normalizar uniformemente sólo el último nivel, ignorando los grados de
       los prefijos y perdiendo la consistencia proyectiva;
 \item usar una lectura posterior para decidir qué estados entran en
       \(\widehat{\mathcal X}^{\rm tot}_k\).
\end{enumerate}
\end{proposition}

\begin{proof}
Los seis primeros puntos eliminan coordenadas o relaciones que aparecen en la
definición de estado, árbol o multisección.  El séptimo elimina la prueba de
\cref{thm:continuo-comun-total-superviviente}.  El octavo puede violar
\eqref{eq:continuo-comun-medida-pushforward} cuando los grados terminales no
son constantes.  El último cambia el dominio total
\eqref{eq:continuo-comun-dominio-total} por una selección retrospectiva.
\end{proof}

Las células APP, la pareja de hojas, el levantamiento TRIT, las extensiones
TPK, el cociclo de acarreo, la memoria, la frontera, la ruta y la conservación
del registro constituyen las estructuras de partida. Sobre ellas se reúnen en
un único dominio el grupoide finito, el árbol de palabras, la multisección y
la medida local normalizada. Los teoremas de supervivencia total, no vacuidad,
consistencia y naturalidad certifican esta construcción conjunta.
